\chapter{The exterior capacitary potential}\label{ch:potential}
\module{NoCompromise.Capacity.Flux,
        NoCompromise.Capacity.Kelvin,
        NoCompromise.Capacity.Levels,
        NoCompromise.Capacity.Potential}

\section{Existence, uniqueness and the Hopf sign}

\begin{theorem}[Capacitary potential]\label{thm:capacitary-potential}
\uses{lem:hull-properties,conv:domain-classes}
There is a unique
\[
  u\in C^\infty(\R^3\setminus K)
    \cap C^{2,\alpha}_{\mathrm{loc}}\bigl(\overline{\R^3\setminus K}\bigr)
  \qquad(0<\alpha<1)
\]
with
\begin{equation}\label{eq:capacitary-problem}
  \Delta u=0\ \text{on }\R^3\setminus K,
  \qquad
  u=1\ \text{on }\partial K,
  \qquad
  u(x)\to0\ (|x|\to\infty).
\end{equation}
It satisfies
\begin{equation}\label{eq:capacitary-signs}
  0<u<1\ \text{outside }K,
  \qquad
  \partial_{\nu_K}u<0\ \text{on }\partial K .
\end{equation}
Extending $u$ by the constant $1$ on $K$ gives a continuous, locally $H^1$
function on $\R^3$.
\end{theorem}

\begin{proof}
\uses{prop:weak-dirichlet,prop:strong-max,lem:weak-max,thm:boundary-C2a,
      cor:hopf-exterior,prop:hopf,lem:sobolev-Ck,prop:poincare-trace,
      lem:hull-properties}
Translate so that $0\in\operatorname{int}K$ and choose
\begin{equation}\label{eq:hull-radii}
  0<a<R_0,\qquad B_a\subset\operatorname{int}K\subset K\subset B_{R_0}.
\end{equation}

\emph{Step 1: annular minimisation.}  For $R>R_0$ let
$A_R:=B_R\setminus K$.  In the affine Sobolev class with traces $1$ on
$\partial K$ and $0$ on $\partial B_R$, minimise
$J(v):=\tfrac12\int_{A_R}|\nabla v|^2$.  A fixed smooth interpolant makes the
class nonempty; \eqref{eq:poincare-L2} on the zero-trace difference space gives
coercivity, and the direct method
(Proposition~\ref{prop:weak-dirichlet}) gives a unique minimiser $u_R$, weakly
harmonic.  Testing with $(u_R-1)_+$ and $(u_R)_-$ gives $0\le u_R\le1$;
Lemma~\ref{lem:sobolev-Ck} and Proposition~\ref{prop:strong-max} give strict
inequalities inside.  Boundary flattening and
Theorem~\ref{thm:boundary-C2a} give $C^{2,\alpha}$ regularity up to either
boundary.  Here $A_R$ is connected: an exterior path may have every portion
outside $B_R$ replaced by a path on a sphere of radius in $(R_0,R)$, which is
disjoint from $K$.

For all sufficiently large $R$, a single smooth function equal to one near
$K$ and zero outside $B_{2R_0}$ is an admissible competitor.  Consequently
$\int_{A_R}|\nabla u_R|^2\le C(K)$ uniformly.

\emph{Step 2: monotone limit.}  If $R_2>R_1$, comparison on $A_{R_1}$ gives
$u_{R_2}\ge u_{R_1}$, since the larger-domain solution is nonnegative on
$\partial B_{R_1}$.  The family is monotone and bounded by one; local elliptic
estimates and monotone convergence give a harmonic limit $u$ with \emph{no
subsequence needed}, and the fixed inner-boundary estimates preserve the trace
$u=1$.  Local weak convergence of the gradients and exhaustion pass the
uniform energy bound to the limit:
\begin{equation}\label{eq:capacity-finite-energy}
 \int_{\R^3\setminus K}|\nabla u|^2<\infty .
\end{equation}

\emph{Step 3: barriers.}  The function $R_0/|x|$ is harmonic on the exterior
(its pole lies in $K$), is $\ge1$ on $\partial K$, and tends to zero, so
comparison gives
\begin{equation}\label{eq:upper-barrier}
  u(x)\le\frac{R_0}{|x|}.
\end{equation}
For a lower barrier on the finite annulus use
\begin{equation}\label{eq:lower-barrier-annulus}
  b_R(r):=\frac{a(R/r-1)}{R-a},
\end{equation}
harmonic for $a<r<R$, equal to $1$ at $r=a$ and $0$ at $r=R$.  Since
$|x|\ge a$ on $\partial K$, $b_R(|x|)\le1=u_R$ there, and both vanish on
$\partial B_R$; comparison gives $u_R\ge b_R$, and letting $R\to\infty$,
\begin{equation}\label{eq:lower-barrier}
  \frac a{|x|}\le u(x)\qquad(|x|\ge R_0).
\end{equation}
So $u\to0$ and $u\not\equiv0$.

\emph{Step 4: uniqueness.}  For two solutions apply
Lemma~\ref{lem:weak-max} to their difference on $B_R\setminus K$; the boundary
difference on $\partial B_R$ tends uniformly to zero.

\emph{Step 5: the Hopf sign.}  Corollary~\ref{cor:hopf-exterior}, using an
exterior tangent ball at $p\in\partial K$ and the radial barrier
\eqref{eq:hopf-barrier}, gives $\partial_{\nu_K}(1-u)>0$.
\end{proof}

\begin{notation}[$w$]\label{not:w}
\uses{thm:capacitary-potential,conv:outward}
$w:=|\nabla u|$.  On $\partial K$,
\begin{equation}\label{eq:w-boundary}
  \nabla u=-w\,\nu_K,\qquad w>0 .
\end{equation}
\end{notation}

\section{Capacity, flux and Kelvin asymptotics}

\begin{definition}[Capacity]\label{def:capacity}
\uses{thm:capacitary-potential}
\begin{equation}\label{eq:capacity}
  \Capa(K):=\frac1{4\pi}\int_{\R^3\setminus K}|\nabla u|^2 ,
  \qquad
  \mathsf C:=\Capa(K).
\end{equation}
\end{definition}

\begin{lemma}[Decay]\label{lem:potential-decay}
\uses{thm:capacitary-potential}
$u(x)=O(|x|^{-1})$ and $|\nabla u(x)|=O(|x|^{-2})$.
\end{lemma}

\begin{proof}\uses{lem:harmonic-deriv,thm:capacitary-potential}
\eqref{eq:upper-barrier} rescaled on dyadic annuli combined with the interior
gradient estimate \eqref{eq:harmonic-deriv}.
\end{proof}

\begin{lemma}[Flux identity]\label{lem:flux-identity}
\uses{def:capacity,not:w,thm:sard-3d}
\begin{equation}\label{eq:flux-boundary}
  4\pi\mathsf C=\int_{\partial K}w\,d\Hs^2 ,
\end{equation}
and for every regular value $s\in(0,1)$,
\begin{equation}\label{eq:flux-level}
  \int_{\{u=s\}}w\,d\Hs^2=4\pi\mathsf C .
\end{equation}
\end{lemma}

\begin{proof}
\uses{thm:W11-gauss-green,lem:potential-decay,thm:sard-3d,not:w}
Integrate by parts on $B_R\setminus K$ using \eqref{eq:W11-GG} and let
$R\to\infty$, using Lemma~\ref{lem:potential-decay}: this gives
\eqref{eq:flux-boundary}.  Applying \eqref{eq:W11-GG} on
$\{s<u<1\}$ instead gives \eqref{eq:flux-level}; regular values are dense by
Theorem~\ref{thm:sard-3d}.
\end{proof}

\begin{lemma}[Kelvin transform and the multipole expansion]\label{lem:kelvin}
\uses{thm:capacitary-potential,prop:kelvin-removable}
The Kelvin transform
\begin{equation}\label{eq:kelvin}
  \widetilde u(z):=|z|^{-1}u\Bigl(\frac z{|z|^2}\Bigr)
\end{equation}
extends smoothly across the origin, so
$\widetilde u(z)=\mathsf C_\infty+a\cdot z+O(|z|^2)$ and, transforming back
with differentiated remainders,
\begin{align}
  u(x)&=\frac{\mathsf C_\infty}{|x|}+\frac{a\cdot x}{|x|^3}+O(|x|^{-3}),
    \label{eq:u-expansion}\\
  \nabla u(x)&=-\mathsf C_\infty\frac{x}{|x|^3}+O(|x|^{-3}),
    \label{eq:grad-u-expansion}\\
  D^2u(x)&=\mathsf C_\infty\Bigl(3\frac{x\otimes x}{|x|^5}-\frac{I}{|x|^3}\Bigr)
    +O(|x|^{-4}).
    \label{eq:hess-u-expansion}
\end{align}
Moreover $\mathsf C_\infty=\mathsf C$.
\end{lemma}

\begin{proof}
\uses{prop:kelvin-removable,lem:potential-decay,lem:flux-identity,
      lem:caccioppoli}
$\widetilde u$ is bounded and harmonic on a punctured ball by
Lemma~\ref{lem:potential-decay}; the removability argument of
Proposition~\ref{prop:kelvin-removable} --- test against a cutoff vanishing on
$B_\varepsilon$, bound the boundary error by $\sup\widetilde u$ times the
capacity of $B_\varepsilon$ --- makes it distributionally harmonic, and
difference-quotient interior regularity makes it smooth.  Taylor's theorem
then gives the expansion, and differentiating the explicit Kelvin
transformation gives \eqref{eq:u-expansion}--\eqref{eq:hess-u-expansion}.  The
inward flux through a large sphere is $4\pi\mathsf C_\infty+o(1)$; comparing
with \eqref{eq:flux-boundary} gives $\mathsf C_\infty=\mathsf C$.
\end{proof}

\begin{lemma}[Asymptotic level geometry]\label{lem:level-asymptotics}
\uses{lem:kelvin,conv:curvature}
For small regular $\varepsilon>0$ the level $\{u=\varepsilon\}$ is a radial
graph $r_\varepsilon(\theta)=\mathsf C/\varepsilon+O(1)$ with two angular
derivatives controlled uniformly, and with the normal toward infinity,
\begin{equation}\label{eq:level-asymptotics}
  w=\frac{\varepsilon^2}{\mathsf C}+O(\varepsilon^3),
  \quad
  H=\frac{2\varepsilon}{\mathsf C}+O(\varepsilon^2),
  \quad
  d\Hs^2=\frac{\mathsf C^2}{\varepsilon^2}\bigl(1+O(\varepsilon)\bigr)d\omega .
\end{equation}
Consequently
\begin{equation}\label{eq:level-Hw-asymptotics}
  \int_{\{u=\varepsilon\}}\bigl(Hw-4\varepsilon^{-1}w^2\bigr)
  =-8\pi\varepsilon+O(\varepsilon^2).
\end{equation}
\end{lemma}

\begin{proof}\uses{lem:kelvin,thm:sard-3d}
The implicit function theorem applied to $u(r\theta)=\varepsilon$ using
\eqref{eq:u-expansion}--\eqref{eq:hess-u-expansion}.
\end{proof}

\begin{fnote}[Off the critical path]\label{fnote:level-asymptotics-unused}
Lemma~\ref{lem:level-asymptotics} is \emph{not used} below.  Chapter~\ref{ch:appK}
needs the far-field normalisation in the translated coordinates of
Lemma~\ref{lem:K-normalization} and re-derives exactly what it needs as
Lemma~\ref{lem:K-far-field}; in particular
\eqref{eq:level-Hw-asymptotics} is the untranslated precursor of
\eqref{eq:K-far-field-limits}.  Lemma~\ref{lem:level-asymptotics} is retained
because it is the natural place to record the level geometry, and because
comparing it with Lemma~\ref{lem:K-far-field} is a useful consistency check
(Regression test~\ref{test:K-model}).  A formalisation may skip it.
\end{fnote}

\section{Connectedness of regular levels}

\begin{lemma}[Regular levels are connected]\label{lem:level-connected}
\uses{thm:capacitary-potential,thm:component-count,lem:hull-properties}
Every regular level $\{u=s\}$, $0<s<1$, is connected.
\end{lemma}

\begin{proof}
\uses{prop:strong-max,lem:weak-max,thm:component-count,lem:hull-properties,
      lem:sobolev-Ck}
Extend $u$ by $1$ on $K$.  The superlevel $\{u>s\}$ is bounded by
\eqref{eq:upper-barrier}.  Every one of its components meets $K$: otherwise
$u$ would be harmonic on a bounded component, equal to $s$ on its boundary and
$>s$ inside, contradicting Proposition~\ref{prop:strong-max}.  Since $K$ is
connected (Lemma~\ref{lem:hull-properties}) it lies in one component, so the
whole superlevel is connected.

The sublevel $\{u<s\}$ contains the complement of a large ball by
\eqref{eq:upper-barrier}.  A bounded component would have boundary value $s$
and an interior minimum below $s$, contradicting the minimum principle.  So the
sublevel is connected.  These are exactly the two complementary components of
the regular level.  That level is a compact smooth surface with finitely many,
say $k$, connected components; Theorem~\ref{thm:component-count} says its
complement has $k+1$ components.  Hence $k+1=2$ and $k=1$.
\end{proof}

